\documentclass[10pt,a4paper]{amsart}
\usepackage[margin=19mm]{geometry}
\usepackage{amsmath,amssymb,mathtools}
\usepackage[scr=rsfs,cal=euler]{mathalfa}
\usepackage{xfrac}
\usepackage[unicode,hidelinks,bookmarks=false]{hyperref}
\newcommand{\ie}{i.e.\ }
\newcommand{\cf}{cf.\ }
\newcommand{\resp}{respectively\ }
\newcommand{\Subsection}{\subsection}
\providecommand{\nobreakdash}{\nobreak\hbox{-}}
\setlength{\emergencystretch}{2em}
\multlinegap0pt
\usepackage{amssymb}
\usepackage{xcolor}
\usepackage{enumerate}
\usepackage{bm}
\usepackage{dsfont}
\usepackage{mathtools}
\newtheorem{theorem}{Theorem}
\title{A skew polynomial framework for constructing division algebras and linear maximum rank distance codes}
\author{S. Pumpl\"un}
\date{Source: arXiv:2606.18371v1 (CC BY 4.0)}
\begin{document}
\maketitle

\noindent\emph{Source and licence.} A skew polynomial framework for constructing division algebras and linear maximum rank distance codes. Version arXiv:2606.18371v1 (CC BY 4.0), distributed by arXiv under the Creative Commons Attribution 4.0 International licence, http://creativecommons.org/licenses/by/4.0/, as recorded in the arXiv licence metadata for this exact version and shown on the abstract page https://arxiv.org/abs/2606.18371v1. Full licence notice: \texttt{licenses/arxiv-2606.18371-CC-BY-4.0.txt}.\par\smallskip

\noindent\textit{Editorial scope.} Counted unit: the article's theorem thm:algebra3 (source0.tex lines 984--996). The article's complete original proof of it is delivered with it (source0.tex lines 998--1031, one \texttt{proof} environment of 34 lines). No mathematics is written, repaired or re-derived: the statement and the proof are the article's own lines, in the article's own notation, and no retained line is substituted or re-flowed.  Retained uncounted as the source-local context the proof needs: lines 933--943.  Nothing was omitted from any retained span, so the package carries no omission record.  No retained line contains a citation, so the wrapper carries no bibliography and no bibliography entry.  No retained line contains a figure environment, an image-inclusion command or drawing code, so no image is delivered and the package declares no asset dependency.  Editorial formatting only, in addition: an AMS \texttt{amsart} preamble that restates the article's own declarations and macros used by the retained lines, namely the environments \texttt{theorem} on separate counters, together with the standard packages those lines need.  The retained environments print in this document as Theorem 1 in order of appearance, whereas the article prints the same items with the numbering of its own full document.  The complete native source of the article as distributed in the arXiv e-print for this exact version is bundled unchanged as \texttt{sources/203075/source0.tex}.
\begin{theorem}\label{thm:algebra2}
Let $f(t)=a_0+a_1t+\dots+ a_{m-1} t^{m-1}+t^m \in R$ be  irreducible  and not two-sided. Let
$$\rho_{\nu,i_0}(y)=y-\nu a_{i_0}\rho(y)$$
 be a bijective $F'$-linear map.
  Let $ g(t)=\sum_{i=0}^{m-1} g_it^i $ and $b(t)=\sum_{i=0}^{m-1} b_it^i$.
 Define
$$g(t) * b(t) =G_{i_0}(H^{-1}(g(t))) b(t){\rm mod}_r\,f$$
$$=\big(\sum_{i=0,i\not=i_0}^{m-1} g_it^i +\rho_{\nu,i_0}(g_{i_0})t^{i_0} + \nu\rho(g_{i_0})f(t)\big) \big( \sum_{i=0}^{m-1} b_it^i\big) {\rm mod}_r\,f,$$
 then $(R_m,*)$ is a non-unital division algebra  of dimension $mn[F:F_0]$ over $F_0$
 and is isotopic to the unital Petit division algebra $\mathbb{S}_f$, viewed as an algebra over $F_0$.
\end{theorem}

\begin{theorem} \label{thm:algebra3}
Let $f\in R$ be monic and irreducible of degree $m$ and not two-sided. Let
$$\rho_{\nu,i_0}(y)=y-\nu a_{i_0}\rho(y)$$
 be a bijective $F'$-linear map.
 Let $z(t^n)\in C(R)$ satisfy that
 $$z(t^n)t^{mk}=1\,{\rm mod}\,f.$$
 Then $R_m$ becomes an $F_0$-algebra  of dimension $mn[F:F_0]$  via the multiplication
$$g(t) \star b(t)=G_{i_0}(H^{-1}(g(t))) \big( z(t^n)t^{mk-1} b(t)\big) \, {\rm mod}_r\, f$$
$$=\big(\sum_{i=0,i\not=i_0}^{m-1} g_it^i +\rho_{\nu,i_0}(g_{i_0})t^{i_0} + \nu\rho(g_{i_0})f(t)\big) z(t^n)t^{mk-1} \sum_{i=0}^{m-1} b_it^i \, {\rm mod}_r\,f$$
when $b(t)=\sum_{i=0}^{m-1} b_it^i$ and $g(t)=\sum_{i=0}^{m-1} g_it^i$.
This algebra is a division algebra and is isotopic to the $F_0$-algebra from Theorem  \ref{thm:algebra2} and to $\mathbb{S}_f$ viewed as an $F_0$-algebra.
\\ (ii) The algebra $(R_m,\star)$  has $t$ as a left unit element, but not as right unit element, and center $F_0$.
\end{theorem}

\begin{proof}
$(i)$ We have
$G_{i_0}(\sum_{i=0,i\not=i_0}^{m-1} g_it^i +\rho_{\nu,i_0}(g_{i_0})t^{i_0})$
$$=\sum_{i=0,i\not=i_0}^{m-1} g_it^i +\rho_{\nu,i_0}(g_{i_0})t^{i_0}+\nu\rho(g_{i_0})f(t)=
\sum_{i=0,i\not=i_0}^{m-1} (g_i +a_i \nu \rho(g_{i_0})  )t^i +   g_{i_0}t^{i_0}+ \nu\rho(g_{i_0})t^m.$$
 Put $H(b(t))= z(t^n)t^{mk-1}b(t)$ then
 $$g(t) \star b(t)=G_{i_0}(H^{-1}(g(t))) F(b(t)) \, {\rm mod}_r\, f$$
 hence $(R_m,\star)$ is isotopic to both $\mathbb{S}_f$ and the algebra from  Theorem  \ref{thm:algebra2}. In particular, it must be a division algebra.
 \\ $(ii)$ Since  we have
 $$g(t) \star t=G_{i_0}(H^{-1}(g(t))) \big( z(t^n)t^{mk}\big) \, {\rm mod}_r\, f$$
$$=\big(\sum_{i=0,i\not=i_0}^{m-1} g_it^i +\rho_{\nu,i_0}(g_{i_0})t^{i_0} + \nu\rho(g_{i_0})f(t)\big) z(t^n)t^{mk} \, {\rm mod}_r\,f$$
$$=\big(\sum_{i=0,i\not=i_0}^{m-1} g_it^i +\rho_{\nu,i_0}(g_{i_0})t^{i_0} + \nu\rho(g_{i_0})f(t)\big)  (s(t)f(t)+1) \, {\rm mod}_r\,f$$
$$=\big(\sum_{i=0,i\not=i_0}^{m-1} g_it^i +\rho_{\nu,i_0}(g_{i_0})t^{i_0} + \nu\rho(g_{i_0})f(t)\big)   \, {\rm mod}_r\,f$$
$$=\big(\sum_{i=0,i\not=i_0}^{m-1} g_it^i +\rho_{\nu,i_0}(g_{i_0})t^{i_0}\big)   \, {\rm mod}_r\,f=g(t)$$
the algebra has a right unit element given by $t$.
Note that for all $i_0\not=1$ we  get
$$ t \star b(t)
=t   z(t^n)t^{mk-1} b(t) {\rm mod}_r\,f=  z(t^n)t^{mk} b(t)\, {\rm mod}_r\,f.$$
Now $ z(t^n)t^{mk}=s(t)f(t)+1$ implies that
$$ t \star b(t)=  (s(t)f(t)+1) b(t) \, {\rm mod}_r\,f$$
$$= (s(t)f(t) b(t) +b(t))\, {\rm mod}_r\,f$$
$$= s(t)f(t) b(t)\, {\rm mod}_r\,f +   b(t)\, {\rm mod}_r\,f.$$
and unless $f$ and $b$ commmute for all $b$, this is generally not equal to $ b(t) \,{\rm mod}_r\,f$, so that $t$ is not a left unit element.
This only is the case when $f\in Z(R)$.

For $i_0=1$ we  get
$$ t \star b(t)
=\big(\rho_{\nu,i_0}(1)t^{i_0} + \nu\rho(1)f(t)\big)  z(t^n)t^{mk-1} b(t) {\rm mod}_r\,f$$
$$= \big(\rho_{\nu,i_0}(1)t^{i_0} (s(t)f(t)+1)  b(t)
+ \nu\rho(1)f(t)  (s(t)f(t)+1)  b(t) \, {\rm mod}_r\,f$$
this is generally not equal to $ b(t) \,{\rm mod}_r\,f$, so that $t$ is not a left unit element.

The algebra $(R_m,\star)$ has center $F_0$ by construction.
\end{proof}

\end{document}
